\chapter{FHCP Challenge y enjambre de partículas}
\section{Métricas del FHCP Challenge}

El Flinders Hamiltonian Cycle Project (FHCP) define el HCP como el problema de determinar si un grafo de \(N\) vértices contiene un ciclo simple de longitud \(N\). \cite{de2}

\subsection{Características del benchmark}

\begin{itemize}
    \item El FHCP Challenge Set contiene 1{,}001 instancias del HCP. \cite{aco9}
    \item Los tamaños de los grafos van desde 66 hasta 9{,}528 vértices. \cite{aco9}
    \item El tamaño promedio es ligeramente superior a 3{,}000 vértices. \cite{aco9}
    \item Las instancias fueron diseñadas para ser difíciles debido a su estructura, no únicamente por su tamaño. \cite{de2}
    \item Todas las instancias tienen al menos un ciclo hamiltoniano conocido. \cite{de2}
\end{itemize}

\subsection{Formato de los archivos}

Los archivos del FHCP utilizan el formato HCP, estructurado en tres secciones:

\begin{itemize}
    \item Encabezado.
    \item Lista de aristas.
    \item Pie de archivo.
\end{itemize}

El campo \texttt{DIMENSION} indica el número de vértices. El campo \texttt{EDGE\_DATA\_FORMAT : EDGE\_LIST} indica que las conexiones se representan mediante una lista de aristas. \cite{aco9}

Los archivos de solución o tour contienen la secuencia de vértices en el orden en que son visitados por el ciclo hamiltoniano. \cite{aco9}

\subsection{Criterios de comparación}

Para realizar una comparación justa entre algoritmos, se deben fijar:

\begin{itemize}
    \item El mismo subconjunto de instancias.
    \item El mismo límite de tiempo o número de evaluaciones.
    \item Varias ejecuciones independientes por instancia.
    \item La misma definición de éxito.
    \item La verificación de que el ciclo encontrado sea válido.
    \item Métricas de tiempo, evaluaciones, tasa de éxito y desviación estándar.
\end{itemize}


\subsection{Reglas históricas y alcance de la comparación}
El concurso histórico estuvo abierto del 30 de septiembre de 2015 al 30 de septiembre de 2016. Ganaba la entrega con más grafos resueltos; en caso de empate se consideraba la recibida primero. El mejor equipo resolvió 985 instancias. Las reglas exigían soluciones en el formato correcto. Esta cobertura es una referencia histórica y no un tiempo de ejecución comparable directamente con nuestro equipo. Las semillas, repeticiones y presupuestos uniformes antes propuestos son un protocolo académico adicional, no exigencias oficiales del concurso.\cite{aco9}

\subsection{Ejemplo completo de entrada y salida}
Este ejemplo propio contiene cinco vértices. Cada par de la entrada representa una arista no dirigida. Los terminadores -1 y EOF no son vértices.
\begin{verbatim}
NAME : ciclo5
COMMENT : Ejemplo didactico
TYPE : HCP
DIMENSION : 5
EDGE_DATA_FORMAT : EDGE_LIST
EDGE_DATA_SECTION
1 2
2 3
3 4
4 5
5 1
-1
EOF
\end{verbatim}
La solución lista cada vértice una sola vez; el cierre es implícito:
\begin{verbatim}
NAME : ciclo5_solucion
COMMENT : Ciclo verificado
TYPE : TOUR
DIMENSION : 5
TOUR_SECTION
1
2
3
4
5
-1
EOF
\end{verbatim}
Se propone validar rangos de identificadores, encabezados y aristas; conservar la numeración de entrada y restaurarla al exportar. El verificador debe comprobar cobertura, unicidad, adyacencias y cierre, no igualdad exacta con un tour de referencia: un grafo puede tener varios ciclos correctos.\cite{aco9}

\section{Particle Swarm Optimization (PSO)}

PSO modela un enjambre de partículas que ajustan su posición según su mejor experiencia personal y la mejor experiencia colectiva.

La actualización de velocidad se define como:

\[
v_i^{t+1}=wv_i^t+c_1r_1(p_i-x_i^t)+c_2r_2(g-x_i^t)
\]

La actualización de posición es:

\[
x_i^{t+1}=x_i^t+v_i^{t+1}
\]

donde:

\begin{itemize}
    \item \(w\) es el coeficiente de inercia.
    \item \(c_1\) y \(c_2\) son coeficientes de aprendizaje.
    \item \(r_1\) y \(r_2\) son valores aleatorios.
    \item \(p_i\) es la mejor posición personal.
    \item \(g\) es la mejor posición global.
\end{itemize}

Para el HCP, PSO requiere una codificación discreta, como random keys, o una variante discreta con operadores definidos sobre intercambios. Las ecuaciones continuas no pueden aplicarse directamente a los identificadores de vértices. \cite{pso1995,pso1998}


\subsection{Alcance de las ecuaciones}
La formulación con peso de inercia $w$ corresponde a la modificación de Shi y Eberhart de 1998; no se atribuye como tal al artículo original de 1995. Los factores aleatorios se generan independientemente por componente y actualización; los productos con vectores se interpretan componente a componente. $p_i$ guarda la posición con mejor costo personal y $g$ la mejor conocida del enjambre en una topología global.\cite{pso1995,pso1998}

\subsection{Decodificación y pseudocódigo propuestos}
Se complementa la propuesta de investigacion.tex con claves aleatorias: una coordenada real por vértice y un ordenamiento para obtener la permutación. Claves $(0.4,0.1,0.8,0.3,0.6)$ para $(A,B,C,D,E)$ producen $(B,D,A,E,C)$. Un empate se resuelve por identificador. La representación garantiza unicidad y cobertura, pero no la existencia de las conexiones. La técnica de claves procede de Bean; usarla con PSO y el costo de conexiones ausentes es aquí una propuesta de adaptación.\cite{de13}
\begin{enumerate}
\item Recibir el grafo, tamaño del enjambre $S$, parámetros $w,c_1,c_2$, semilla y presupuesto. Definir límites: por ejemplo, recortar claves a $[0,1]$ y velocidades a $[-v_{\max},v_{\max}]$.
\item Inicializar posiciones y velocidades, ordenar claves y evaluar el costo $b$ de cada permutación. Contar también estas evaluaciones. Guardar los mejores personales y global y verificar cualquier costo cero.
\item Mientras queden recursos, actualizar velocidad y posición de cada partícula, y aplicar los límites establecidos.
\item Decodificar, evaluar y contabilizar el esfuerzo. Actualizar el mejor personal si el costo disminuye; conservar el anterior en empate.
\item Actualizar el mejor global. En esta propuesta se hace después de cada partícula; una variante sincrónica lo haría al final de la iteración.
\item Verificar cualquier candidata de costo cero y devolverla si es válida. Si se agotan los recursos, devolver la mejor candidata con el estado ``no encontrado'', sin afirmar inexistencia.
\end{enumerate}

\subsection{Ejemplo didáctico y dificultades}
Sobre $E=\{AB,BC,CD,DE,EA\}$, las claves $(0.1,0.2,0.4,0.3,0.5)$ producen $(A,B,D,C,E)$, con dos conexiones ausentes. Un desplazamiento $(0,0,-0.1,0.1,0)$ produce claves ordenadas y la ruta $(A,B,C,D,E)$: se verifican $AB,BC,CD,DE,EA$. Ese desplazamiento puede corresponder a una velocidad previa igual, con $w=1$ y términos de atracción nulos en ese paso. Es una ilustración algebraica elegida a mano, no una ejecución ni una recomendación de parámetros.

La decodificación cuesta $O(n\log n)$ con ordenamiento por comparación; verificar conexiones cuesta $O(n)$ con consultas constantes de adyacencia. Almacenar el enjambre y sus recuerdos requiere $O(Sn)$ espacio además del grafo. Diferentes vectores pueden decodificar a la misma permutación, produciendo mesetas. La convergencia del enjambre puede reducir diversidad sin alcanzar un ciclo. Reinicios, límites y topología social son decisiones que deben documentarse y medirse, sin prometer éxito en un presupuesto finito.
